%---------------------------------------------------------------------
\chapter{Phase 2 --- A Fail-Closed Conformance Layer}
\label{ch:conformance}

Phase 1 ended with the observation that the \emph{measurement} of zero-variance
structure is reliable while most \emph{algorithmic} levers tested were
noise-limited. Phase 2 therefore did not chase another GRPO variant. It built the
layer the Phase-1 synthesis pointed at: an auditable chain from the stated
objective to the recorded training decision, and a benchmark-independent standard
for reporting what a run actually did. This chapter specifies that layer. It is
the methodological spine on which the fourteen-lane campaign of
Chapters~\ref{ch:campaign}--\ref{ch:p2results} rests.

\section{Motivation: the same terminal signal can demand opposite actions}

Group-relative RL can assign exactly zero centered advantage to a group whose
rewards are all equal --- the $\zvf$ event Phase 1 measured. The event is trivial
to detect and entirely ambiguous to interpret. An all-zero group may reflect:

\begin{itemize}[leftmargin=1.4em,itemsep=2pt]
  \item an \emph{impossible} prompt, where no completion can succeed and the group
  should be dropped;
  \item a \emph{capable} policy that was unlucky on this sample, where the group is
  informative and dropping it discards signal;
  \item a \emph{false-negative verifier}, where the reward channel is broken and
  the correct action is to halt rather than adapt;
  \item an \emph{execution failure}, where the infrastructure, not the policy, is at
  fault.
\end{itemize}

These four states can require \textbf{opposite} decisions --- resample, retain,
halt, or retry --- despite exposing an identical terminal reward history. A
controller that keys only on ``was the group zero-variance?'' is therefore not
merely imprecise; it is decision-insufficient. This is the ambiguity that
motivates a conformance layer rather than another sampling heuristic.

\section{Reference semantics and the frozen differential}

The harness evaluates frozen reward and mask fixtures on CPU without loading a
model or running an optimizer. A framework-neutral reference emits advantages,
selected rows, losses, and decision outputs. Separate adapters evaluate pinned
TRL and verl distributions and record package versions, source paths, and
hashes. The combined receipt carries a SHA-256 fixture digest and contains
\textbf{14 intended cases per stack} plus a \textbf{shared 36-case controller
matrix}, with no recorded errors.

\subsection{The scalar objective under test}

Let $\mathcal{S}$ be the selected completion rows and $M_{it}\in\{0,1\}$ the
completion mask. Within each reward group, GRPO, DAPO, and GSPO use the sample
standard deviation (Bessel correction); if it is at most $10^{-8}$, $A_i=0$.
DrGRPO uses only the centered reward. AERO replaces the zero-variance branch with
its frozen Beta-smoothed posterior scale. For token log probabilities
$\ell_{it}$ and old log probabilities $\ell^{\mathrm{old}}_{it}$, the exact
advantage rule is

\begin{equation}
A_i=\begin{cases}
(R_i-\bar R_g)/s_g, & \text{arm}\ne\text{DrGRPO},\ s_g>10^{-8},\\[2pt]
R_i-\bar R_g, & \text{arm}=\text{DrGRPO},\\[2pt]
(R_i-\tilde p_g)/\sqrt{\tilde p_g(1-\tilde p_g)},
  & \text{arm}=\text{AERO},\ s_g\leq10^{-8},\\[2pt]
0, & \text{otherwise},
\end{cases}
\quad \tilde p_g=\frac{k_g+1}{n_g+2},
\label{eq:s1-advantage}
\end{equation}

where $\bar R_g$ and $s_g$ use the active rows of group $g$, and $(k_g,n_g)$ are
the frozen AERO success/observation counts. The loss is

\begin{align}
 \rho_{it}&=\exp(\ell_{it}-\ell^{\mathrm{old}}_{it}), \nonumber\\
 h_{it}&=-\min\!\left\{\rho_{it}A_i,
     \operatorname{clip}(\rho_{it};0.8,1+u)A_i\right\}, \nonumber\\
 L_{\mathrm{ref}}&=\frac{1}{|\mathcal{S}|}\sum_{i\in\mathcal{S}}
   \frac{\sum_t M_{it}h_{it}}{\sum_t M_{it}},
 \label{eq:s1-objective}
\end{align}

where $u=0.28$ for DAPO and $0.20$ otherwise. GSPO replaces the token ratio by
the sequence ratio
$\exp(\sum_t M_{it}(\ell_{it}-\ell^{\mathrm{old}}_{it})/\sum_t M_{it})$. There
is no KL or other auxiliary term in this CPU differential.

\subsection{Four campaign conditions}

The four campaign conditions use the DAPO loss structure of
Eq.~\eqref{eq:s1-objective}, but the training implementation carries a separate
$10^{-12}$ zero-standard-deviation floor (S1 uses $10^{-8}$ for its float64
conformance comparisons):

\begin{itemize}[leftmargin=1.4em,itemsep=2pt]
  \item \texttt{intended\_full} --- standardizes selected rewards without an
  additive epsilon, emits zero advantages at the campaign floor, and averages
  completion means.
  \item \texttt{epsilon\_only} --- adds $10^{-4}$ to a nonzero selected standard
  deviation.
  \item \texttt{reduction\_only} --- keeps the intended advantages but uses one
  global masked-token mean.
  \item \texttt{native\_trl} --- standardizes all eight reward rows with
  $s+10^{-4}$ \emph{before} selection and uses the global masked-token mean; it
  does not use the campaign's explicit floor branch.
\end{itemize}

Inactive rows carry all-zero masks. The verifier reads and asserts both source
constants, so a silent change to either floor fails the protocol rather than
altering the verdict.

The harness deliberately separates two questions:

\begin{enumerate}[leftmargin=1.6em,itemsep=3pt]
  \item Does the intended integration reproduce the canonical reference within
  absolute tolerance $10^{-8}$ and relative tolerance $10^{-6}$ in float64?
  \item What does the stack's \emph{native} objective do on the same fixture?
\end{enumerate}

The combined receipt is \texttt{S1\_PASS} for the first question. For the second,
it preserves native verdicts without relabelling them as conformance failures.
Table~\ref{tab:s1} records the split. These verdicts mean ``different from this
reference,'' not ``framework bug.''

\begin{table}[h]
  \centering
  \caption{Frozen S1 conformance receipt. ``Material'' and ``not tested'' are
  native-versus-reference verdicts; intended integrations pass separately.}
  \label{tab:s1}
  \begin{tabular}{lrrr}
    \toprule
    Stack & Intended cases & Native material & Native not tested \\
    \midrule
    TRL 1.2.0 & 14 & 4 & 1 \\
    verl 0.3.0.post1 & 14 & 1 & 4 \\
    \bottomrule
  \end{tabular}
\end{table}

\subsection{The controller matrix}

The controller matrix is shared across stacks: six policies crossed with six
observations. The policies are fixed $G\in\{8,16\}$, symmetric-ZVF,
failure-only, boundary-aware, and full-triage. The observations are all-wrong,
all-correct, mixed, noisy, missing, and delayed. Every policy returns
\texttt{recheck} for an unresolved reward. Otherwise:

\begin{itemize}[leftmargin=1.4em,itemsep=2pt]
  \item static policies keep their named size;
  \item \textbf{symmetric-ZVF} escalates any zero-variance group to 16;
  \item \textbf{failure-only} escalates all-wrong groups;
  \item \textbf{boundary-aware} escalates all-wrong and drops all-correct groups;
  \item \textbf{full-triage} escalates all-wrong when the Wilson upper bound is at
  least $0.05$ (dropping it otherwise), drops all-correct when the Wilson lower
  bound is at least $0.95$, and keeps $G=8$ otherwise.
\end{itemize}

\section{Zero gradients are relations, not angles}

Let $g_a$ and $g_b$ be gradients under two semantics. A cosine is meaningful only
when both norms are positive. Here each gradient is the float64 CPU gradient of
the complete selected policy objective in Eq.~\eqref{eq:s1-objective}; ``zero''
means a bit-exact stored norm of zero, not a tunable threshold. Receipts retain
both norms and use four logical cases:

\begin{equation}
\operatorname{rel}(g_a,g_b)=
\begin{cases}
  \texttt{joint\_zero}, & \lVert g_a\rVert=\lVert g_b\rVert=0,\\
  \texttt{a\_zero}, & \lVert g_a\rVert=0<\lVert g_b\rVert,\\
  \texttt{b\_zero}, & \lVert g_b\rVert=0<\lVert g_a\rVert,\\
  \texttt{nonzero}, & \lVert g_a\rVert\lVert g_b\rVert>0.
\end{cases}
\label{eq:gradient-relation}
\end{equation}

Cosine and relative $\ell_2$ distance are null in the first three cases. Joint
zero counts as equivalence; a one-sided zero is a categorical material
divergence. This classification alone makes no general claim about parameter
updates under momentum, weight decay, auxiliary losses, or accumulation. In the
frozen campaign specifically, $\beta=0$, weight decay is zero, gradient
accumulation is one, and a selected zero gradient explicitly skips
\texttt{optimizer.step}; parameters and AdamW state are unchanged while the
linear scheduler advances. This avoids both a division-by-zero artifact and the
selective deletion of hard replay groups.

\section{Cost-accounting corollaries of the starvation model}

Phase 1 modelled starvation descriptively. Phase 2 added cost-accounting
corollaries that hold for the known binary-starvation model, giving the
controller a principled reason to prefer some escalations over others:

\begin{itemize}[leftmargin=1.4em,itemsep=2pt]
  \item \textbf{Contrast is monotone in $G$.} For fixed $G$, increasing group size
  reduces the expected number of groups needed to observe a contrast.
  \item \textbf{Complete-group sampling is not free.} For every $G$,
  complete-group sampling costs at least $1/\epsilon$ expected rollouts to reach
  success probability $\epsilon$ on a group whose success probability is bounded
  below by $\epsilon$.
  \item \textbf{Decision reversal.} A restricted two-state witness shows that when
  $\Delta>0$, the two actions $\{\text{stop},\ \text{commit to }m\text{ batch
  rollouts}\}$ reverse rank across compatible latent states and incur positive
  two-point minimax regret.
\end{itemize}

The practical consequence is that ``resample until a contrast appears'' is not a
free repair. It has a quantified rollout cost, and under latent-state uncertainty
the better action can reverse. That is precisely why the campaign's lanes are
budgeted and gated rather than open-ended.

\section{Failure boundaries and immutable evidence}

The protocol \textbf{rejects} a unit when any required source hash, corpus
fingerprint, package pin, profiler phase, gradient relation, checkpoint manifest,
or matched-budget field is absent or inconsistent. Three design commitments
follow, and they govern every result reported in Chapters~\ref{ch:campaign} and
\ref{ch:p2results}:

\begin{enumerate}[leftmargin=1.6em,itemsep=3pt]
  \item \textbf{Infrastructure attempts are not scientific units.} Failed
  attempts remain separate from accepted units; they are never pooled into a
  denominator or averaged into a score.
  \item \textbf{Costs are attributed, not blended.} Corpus generation, diagnostic
  backward passes, and held-out evaluation are charged explicitly, so a lane's
  spend is auditable to its receipts.
  \item \textbf{Retries are constrained.} A retry may resume only the greatest
  checkpoint whose manifest and content fingerprints match.
\end{enumerate}

Development logs record several infrastructure and diagnostic corrections, but
this report does not treat those anecdotes as scientific evidence. The claim
boundary begins at immutable accepted receipts. In particular, a failed
positive-control construction is reported as a \emph{registered feasibility
outcome}, not quietly rerun and replaced with a passing version.

\section{The registered feasibility outcome}

The frozen screening campaign reported in this chapter produced the following
outcome, stated without filling in missing evidence:

\begin{itemize}[leftmargin=1.4em,itemsep=2pt]
  \item \textbf{Six scientific units completed} in the balanced equal-length
  negative control.
  \item Each contained \textbf{100 replay steps}, of which \textbf{62--65 store
  exact joint-zero gradient diagnostics}.
  \item All remaining intended-versus-native cosines were at least
  \textbf{0.999844}.
  \item High cosine was \emph{not} the registered gate: the sole completed
  intended-full cell satisfies the joint cosine/relative-$\ell_2$ criterion on
  \textbf{69/100 steps}, below the required \textbf{95/100}.
  \item The filtered variable-length positive control was
  \textbf{contract-infeasible}; four seed-37 units remain pending for accelerator
  quota; and \textbf{no confirmatory matrix ran}.
\end{itemize}

Consequently this registered feasibility outcome establishes \emph{neither} a
causal training effect \emph{nor} a scientific negative result about objective
semantics. It establishes what the protocol was able to prove on the compute it
had. The contribution of this chapter is the auditable boundary between
conformance evidence, mathematical necessity, and a registered feasibility
outcome --- not a headline number.
